%% conclusion_ko.tex
\chapter{결론}
\label{chap:conclusion_ko}

\section{연구 요약}
\label{sec:conc_summary_ko}

본 논문은 기업 지식 관리 시스템에서 정보 검색을 향상시키는 메커니즘으로서의 적응형 검색 계획을 연구하였다. 검색 아키텍처 설계, 적응형 계획, 계획기 비교, GPT-5를 활용한 LLM 기반 계획에 걸친 7개의 연구 질문을 다루었다. 모든 실험은 6개 기업 문서 범주와 6가지 추론 유형에 걸친 405개 질의를 포함하는 완전한 주석 120문서, 1,206청크 코퍼스인 합성 기업 지식 데이터셋(SEKD)에서 수행되었다.

실험 프로그램은 9개의 단계를 거쳤다:

\begin{itemize}
    \item \textbf{5단계 (BM25)}: Recall@10$={0.713}$, MRR$={0.457}$의 희소 검색 기준선 확립. 예외 처리 및 시간적 질의에서의 지속적 취약점 규명.
    \item \textbf{6단계 (Dense)}: BGE-small-en-v1.5 밀집 검색이 MRR을 +9.6\%\, 향상시켜 $0.501$을 달성함을 입증. 의미적 추론 과제에서 특히 향상. 단, 시간적 질의에서 리콜 0\%.
    \item \textbf{7단계 (Hybrid)}: Hybrid RRF를 최선의 고정 전략 기준선으로 확립 (Recall@10$={0.745}$, MRR$={0.530}$).
    \item \textbf{8단계 (규칙 계획기)}: 7개 규칙 결정론적 계획기가 SSA$={78.4\%}$, Recall@10$={0.755}$를 달성하여 운영 비용 없이 커버리지 지표에서 Hybrid를 능가함을 입증.
    \item \textbf{9단계 (LLM 계획기)}: GPT-5를 제로샷 검색 전략 계획기로 평가. SSA$={77.6\%}$, MRR$={0.533}$, nDCG@10$={0.570}$ 달성. 순위 품질에서 규칙 계획기에 근소하게 우위, 커버리지에서 근소하게 열세. 380개 질의 처리 비용 \$6.18.
\end{itemize}

\section{연구 질문 답변}
\label{sec:conc_rqs_ko}

\subsection*{RQ1: 하이브리드 검색은 단일 방법 기준선을 능가하는가?}
\textbf{그렇다.} Hybrid RRF는 Recall@10$={0.745}$, MRR$={0.530}$, nDCG@10$={0.566}$을 달성하여 BM25(Recall@10$={0.713}$, MRR$={0.457}$)와 Dense(Recall@10$={0.740}$, MRR$={0.501}$) 모두를 1차 평가 지표에서 능가한다. BM25 대비 MRR +15.8\%, Dense 대비 MRR +5.6\%의 향상은 기업 지식 검색에서 어휘적 신호와 의미적 신호가 상보적임을 확인한다.

\subsection*{RQ2: 다양한 범주에서 어떤 검색 방법이 최선인가?}
\textbf{어떤 단일 방법도 모든 범주를 지배하지 않는다.} Hybrid는 6개 범주 전반에 걸쳐 가장 균형 잡힌 성능을 제공한다. 시스템 설계 문서와 보안 정책은 Dense를 강하게 선호한다(BM25 대비 Recall@10 각각 +18.9\%, +11.7\%). 출장 정책은 모든 방법에서 완벽한 Recall@10을 달성한다. 인사 정책은 모든 방법에서 가장 약한 범주이다(Recall@10 $={0.546}$--$0.611$).

\subsection*{RQ3: 추론 유형별로 성능은 어떻게 달라지는가?}
\textbf{실질적으로 달라진다.} 집계 및 비교 질의는 모든 시스템에서 Recall@10 $>{0.818}$을 달성한다. 단일 홉 질의는 적당히 양호하다($0.764$--$0.832$). 다중 홉 질의는 어렵다($0.491$--$0.538$). 예외 및 시간적 질의는 모두 공통적으로 어려우며(각각 Recall@10 $\leq{0.412}$와 $\leq{0.200}$), Dense는 시간적 질의에서 리콜 0\%를 기록한다. 최선과 최악의 경우 추론 유형 간 50+ pp 성능 범위는 기업 RAG 시스템 설계에서 질의 유형이 검색 난이도의 지배적 예측 변수임을 확인한다.

\subsection*{RQ4: 규칙 기반 계획기가 고정 전략 검색을 능가할 수 있는가?}
\textbf{그렇다, 근소하게.} 규칙 계획기는 Recall@10$={0.755}$(Hybrid 대비 +1.4\%)와 Hit@10$={0.808}$(Hybrid 대비 +1.3\%)을 달성하며, 대부분의 순위 지표에서 비슷한 성능을 보이고 운영 비용 없이 순 커버리지 향상을 제공한다. 이는 기업 RAG 시스템의 실용적 향상으로서 적응형 계획의 유효성을 확인한다.

\subsection*{RQ5: LLM 기반 계획에 최적인 프롬프트 전략은 무엇인가?}
\textbf{구조화 프롬프트이다.} zero\_shot(SSA$={69.5\%}$), structured(SSA$={82.1\%}$), few\_shot(SSA$={81.8\%}$) 중 구조화 프롬프트가 few\_shot의 약 절반 토큰 비용으로 최고의 SSA를 달성한다. 추론 유형, 난이도 요인, 문서 범주 등의 명시적 메타데이터 제공이 문맥 내 예시 없이도 전략 선택 논리 적용을 가능하게 한다.

\subsection*{RQ6: GPT-5는 전략 선택 정확도에서 규칙 계획기를 능가하는가?}
\textbf{아니다.} GPT-5의 SSA $={77.6\%}$는 규칙 계획기의 $78.4\%$보다 $-$0.79 pp 낮다. GPT-5는 순위 품질 지표에서 규칙 계획기를 능가하지만(MRR: +2.0\%, nDCG@10: +1.2\%), 원시 커버리지에서는 열세이다(Recall@10: $-$0.7\%). 성능 격차는 oracle 레이블 모호성에서 예상되는 범위 내에 있다(GPT-5 오류의 44.7\%가 \texttt{oracle\_disagreement}로 분류됨).

\subsection*{RQ7: LLM 기반 계획의 주요 실패 유형은 무엇인가?}
\textbf{범주 혼동과 시간적 오라우팅이다.} GPT-5의 지배적 실패 유형은 (1) API 문서 및 시스템 설계 질의에서의 범주 혼동(오류의 55.3\%), bm25, dense, hybrid 구분이 모호한 경우; (2) 시간적 질의 실패(SSA 0\%, hybrid를 체계적으로 선호)이다. 단일 홉 질의가 GPT-5 계획 오류의 94\%를 차지하며, LLM 추론이 잘 구조화된 메타데이터 풍부 질의에서 결정론적 규칙에 비해 거의 이점을 제공하지 않음을 보여준다.

\section{주요 연구 발견}
\label{sec:conc_findings_ko}

V1 실험 프로그램에서 세 가지 주요 발견이 도출된다.

\textbf{발견 1: 검색 아키텍처가 성능의 1차 동인이다.} BM25에서 Hybrid RRF로의 향상(MRR +15.8\%)은 최선의 계획기 향상(MRR +2.0\%)보다 8배 크다. 기업 RAG 실무자는 계획 지능보다 검색 아키텍처 선택을 1차 설계 결정으로 우선시해야 한다. 계획 지능 없는 잘 구성된 Hybrid RRF 기준선이 BM25 전용 백엔드 위의 정교한 계획기를 크게 능가한다.

\textbf{발견 2: 결정론적 계획이 수백만 배 낮은 비용으로 LLM 계획과 경쟁적이다.} 7개 규칙 계획기와 GPT-5는 단말 검색 품질에서 기능적으로 동등하다(Recall@10 0.7 pp 이내, MRR 1.1 pp 이내). 처리량 요건이나 비용 제약이 있는 기업 배포에서 규칙 계획기가 우월한 선택이다. 정밀도가 중요한 저용량 응용에서는 GPT-5의 순위 품질 이점이 비용을 정당화할 수 있다.

\textbf{발견 3: 예외 처리 및 시간적 추론이 보편적 실패 유형이다.} 모든 V1 시스템은 예외 조항 질의(Recall@10 $\leq{0.412}$)와 시간적 추론 질의(Recall@10 $\leq{0.200}$)에서 체계적으로 실패한다. 이러한 실패 유형은 검색 전략 선택만으로는 해결할 수 없다. V2 에이전틱 RAG 확장(제~\ref{chap:v2_methodology_ko}--\ref{chap:v2_results_ko}장)은 검증 신호에 조건화된 반복적 재검색을 통해 이러한 실패 모드를 직접 해결한다.

\section{연구 기여}
\label{sec:conc_contributions_ko}

\textbf{방법론적 기여}: SEKD 코퍼스는 섹션 인식 청크 분할, 6개 추론 유형, 8개 난이도 요인, oracle 전략 레이블을 갖춘 완전한 주석 기업 검색 평가 환경을 제공하며, 재현 가능한 실험을 지원한다.

\textbf{실증적 기여}: BM25, 밀집 검색, Hybrid RRF, 규칙 기반 계획기, LLM 기반 계획기를 통제된 기업 지식 코퍼스에서 최초로 체계적으로 비교한 결과를 제공한다.

\textbf{실용적 기여}: 7개 규칙 결정론적 계획기가 \$0의 운영 비용으로 모든 평가 시스템 중 최고의 Recall@10(0.755)을 달성함을 입증하여, 기업 RAG 아키텍트에게 구체적이고 구현 가능한 계획 메커니즘을 제공한다.

\textbf{분석적 기여}: GPT-5의 주요 실패 유형(범주 혼동, 시간적 오라우팅)을 규명하고, 모든 검색 방법에 걸친 예외 처리 및 시간적 질의 실패의 구조적 원인을 식별하였다.

\section{V2 에이전틱 RAG: 요약}
\label{sec:conc_v2_summary_ko}

모든 V1 시스템에서 예외 및 시간적 질의의 지속적 실패에 동기화되어, 제~\ref{chap:v2_methodology_ko}장과 제~\ref{chap:v2_results_ko}장은 V2 에이전틱 RAG 파이프라인을 소개하고 평가한다. 핵심 아키텍처 혁신은 5개 에이전트의 LangGraph 상태 기계로, 검증 에이전트가 각 검색 시도 후 검색 청크 관련성을 평가하고, 검색 세트가 유형별 관련성 임계값을 충족하지 못할 때 타겟 재검색을 트리거한다. 6개의 V2 연구 질문은 다음의 점진적 기여를 다룬다: (RQ8) V1 계획기 대비 완전한 에이전틱 파이프라인; (RQ9) 단일 패스 에이전틱 기준선 대비 재검색 루프; (RQ10) 순위 순서 스코어링 대비 적응형 검증 스코어링; (RQ11) Oracle vs. 휴리스틱 질의 분석; (RQ12) 예외 및 시간적 실패를 특정적으로 해결하는 파이프라인 능력; (RQ13) 루프 소진 후 남은 실패 모드.

\section{향후 연구 방향}
\label{sec:conc_future_ko}

\textbf{예외 및 시간적 검색 전문화.} 가장 중요한 향상 기회는 예외 및 시간적 질의에서 0.000--0.412 Recall@10 성능 하한을 해결하는 것이다. 구체적 방향으로는 시간적 질의를 위한 BM25의 날짜 필드 부스팅, 정책 예외 조항 쌍에 대한 밀집 인코더 미세조정, 날짜 토큰 질의에 BM25를 적용하면서 의미적 예외 질의에는 밀집 검색을 보존하는 하이브리드 전략 개발이 있다.

\textbf{경험적 도출 Oracle 레이블.} 휴리스틱 oracle 레이블을 경험적 도출 레이블(보류된 검색 실험에서 각 질의당 가장 높은 Recall@10을 달성하는 전략)로 대체하면 더 객관적인 SSA 평가 기준선을 제공할 것이다. 이는 현재 \texttt{oracle\_disagreement}로 분류된 44.7\% GPT-5 오류도 명확히 할 것이다.

\textbf{미세조정 검색 계획기.} GPT-5 계획 결정을 레이블로 사용하여 훈련된 소규모 미세조정 분류기(예: 보류된 인코더)는 규칙 계획기의 비용 효율성과 GPT-5에서 관찰된 보정 품질 이점을 결합할 수 있다. 380개 레이블 계획 결정의 가용성을 고려하면 이 접근법은 무시할 수 있는 추론 비용으로 두 기준선 모두를 능가할 가능성이 있다.

\textbf{실제 기업 코퍼스 검증.} SEKD 발견은 언어적 다양성, 문서 간 종속성, 질의 자연성의 영향을 평가하기 위해 실제 기업 데이터셋에서 검증되어야 한다.

\textbf{V2 응답 생성 품질.} V2 응답 에이전트는 인라인 인용 마커가 있는 응답을 합성한다. 향후 연구는 자동 지표(BERTScore, ROUGE-L)와 인간 평가자를 사용하여 응답 품질을 평가하고, 인용과 응답 텍스트 간의 사실적 일관성을 측정해야 한다. 이는 검색 품질에서 하류 응답 정확도로의 루프를 완성한다.

\textbf{재검색 루프의 질의 확장.} 현재 V2 재검색 루프는 검색 전략을 변경하지만 질의를 다시 작성하지 않는다. 검증 에이전트의 실패 진단에서 제안된 질의 확장 힌트를 재발행된 하위 질의에 통합하는 것은 자연스러운 다음 단계로, 특히 원시 질의 텍스트에 관련 문서에 존재하는 정확한 어휘가 없는 \texttt{missing\_entity} 실패에 효과적이다.

\section{맺음말}
\label{sec:conc_closing_ko}

본 연구는 적응형 검색 계획이 고정 전략 기준선을 넘어 기업 정보 검색을 향상시키지만, 계획 메커니즘의 선택이 검색 아키텍처의 선택보다 중요성이 낮음을 입증하였다. Hybrid RRF가 권장 검색 기반이며, 경량 규칙 기반 계획기가 실제 기업 RAG 배포에서 최선의 비용-성능 상충관계를 제공한다.

GPT-5가 7개 규칙 결정론적 계획기를 명확히 능가하거나 명확히 열세하지 않는다는 발견은 부정적 결과가 아니라 현재 기술 최전선의 정보적 특성화이다. 이는 구조화된 메타데이터(추론 유형, 난이도 요인, 문서 범주)에서 가용한 정보가 대부분의 기업 질의에 대한 최적 검색 전략을 결정하기에 충분하며, LLM 추론이 현재 형태의 이 특정 과제에서 한계적 가치를 추가함을 시사한다. LLM 추론 비용이 하락하고 모델 보정이 향상됨에 따라 비용-편익 계산이 GPT-5 방향으로 이동할 것이다. 현재로서는 결정론적 계획이 기업 RAG 배포를 위한 실용적 선택이다.
